\section{Evolutionary History of High-complexity Case Study}
\label{sec:history}

% https://mybinder.org/v2/gh/mmore500/dishtiny/e17e6d5e258b7aacac72d44922008ab14e80e182?filepath=binder%2Fbucket%3Dprq49%2Fa%3Dall_stints_all_series_profiles%2Bendeavor%3D16%2Fcase_study_16005.ipynb
Due to the parallel, asynchronous nature of the experimental framework, we did not perform perfect phylogeny tracking \citep{moreno2024analysis}.
\dissertationonly{Chapter \ref{ch:distributed-phylogeny} discusses challenges parallelizing perfect phylogeny tracking in depth.}

Instead, we opted to track the total number of ancestors seeded into stint 0 with extant descendants.
At the end of stints 0 and 1, three distinct original phylogenetic roots were present in the population.
From stint 2 onward, only two distinct original phylogenetic roots were present.

% https://mybinder.org/v2/gh/mmore500/dishtiny/e17e6d5e258b7aacac72d44922008ab14e80e182?filepath=binder%2Fbucket%3Dprq49%2Fa%3Dall_stints_all_series_profiles%2Bendeavor%3D16%2Fcase_study_16005.ipynb
We performed follow-up analyses on specimens sampled from the lowest original phylogenetic root ID present in the population.%
\footnote{
This approach was designed to choose an arbitrary strain as focal.
Barring extinction, that same strain will then be identified as focal consistently across subsequent stints.
Phylogenetic root ID had no functional consequences; it is simply an arbitrary basis for focal strain selection.
}
For the first two stints, the focal strain was root ID 2,378.
During stint 2, original phylogenetic root 2,378 went extinct.
So, all further follow-up analyses were sampled from descendants of ancestor 12,634.

We also tracked the number of genomes reconstituted at the outset of each stint that had extant descendants at the end of that stint.
This count grows from approximately 10 around stint 15 to upwards of 30 around stint 40 (Supplementary Figure \ref{fig:phylogeny:stint_roots}).
Among descendants of the lowest original phylogenetic root, the number of independent lineages spanning a stint also increases from around 5 to around 15
(Supplementary Figure \ref{fig:phylogeny:lowestroot_stint_roots}).
This decrease in phylogenetic consolidation on a stint-by-stint basis correlates with the waning number of simulation updates performed per stint (Supplementary Figures \ref{fig:phylogeny:updates_vs_stint_roots} and \ref{fig:phylogeny:log_updates_vs_stint_roots}).
More complete phylogenetic data will be necessary in future experiments to address questions about the possibility of long-term stable coexistence beyond the two strains supported under the explicit diversity maintenance scheme.

% https://mybinder.org/v2/gh/mmore500/dishtiny/e17e6d5e258b7aacac72d44922008ab14e80e182?filepath=binder%2Fbucket%3Dprq49%2Fa%3Dall_stints_all_series_profiles%2Bendeavor%3D16%2Fcase_study_16005.ipynb
On the specimen from stint 100 used in the final case study, an evolutionary history of 20,212 cell generations had elapsed.
Of these cellular reproductions, 11,713 (58\%) had full kin group commonality, 7,174 had partial kin group commonality (35\%), and 1,325 had no kin group commonality (7\%).
On this specimen, 1,672 mutation events had elapsed.
During these events, 7,240 insertion-deletion alterations had occurred and 26,153 point mutations had occurred.
This strain experienced a selection pressure of 18\% over its evolutionary history, meaning that only 82\% of the mutations that would be expected given the number of cellular reproductions that had elapsed were present.

% https://github.com/mmore500/dishtiny_event_tag_phylogenetics/commit/ecc2ab8a580b125bea9f661e49cd4d7f34bfa1bb
In order to characterize the evolutionary history of the experiment in greater detail, we performed a parsimony-based phylogenetic reconstruction on the sampled representative specimens from each stint, shown in Supplementary Figure \ref{fig:phylo_parsimony_tree}.
We used genomes' fixed-length blocks of 35 64-bit tags that mediate environmental interactions as the basis for this reconstruction.
These tag blocks underwent bitwise mutation over the course of the experiment.%
\footnote{
In future experiments, we plan to incorporate new methodology for ``hereditary stratigraph'' genome annotations expressly designed to facilitate phylogenetic reconstruction \dissertationelse{(Chapter \ref{ch:distributed-phylogeny})}{\citep{moreno2022hereditary}}.
}
Supplementary Figure \ref{fig:phylo_distance_matrix_heatmap} shows Hamming distance between all pairs of tag blocks.
We additionally tried several other tree inference methods, discussed in supplementary material; however, these yielded lower-quality reconstructions.

Although the phylogeny of stint representatives includes many instances that do not constitute a strict sequential lineage (i.e., each stint's representative descending directly from the preceding stint's representative), we did not observe evidence of long-term coexistence of clades over more than ten stints.

% \subsection{Multicellular Phenotypic Traits}

% In order for a transition in individuality, cells have to diminish their own immediate reproductive interests in order to prioritize the interests of the collective.
% In previous work with an earlier version of the system, we've seen multicellular traits evolve such as (blah from abstract) %todo
% \citep{moreno2021exploring}.

% The formation of kin group morphologies in this system does not necessarily imply a transition in individuality.
% It is necessary to test to see if cells within groups are meaningfully cooperating.

% Because cell reproduction is necessarily adversarial in this system, we can test to see whether cells preferentially target or spare their group members.
% A conflict ratio of less than 1 implies that cells preferentially spare their members.
% We observed a  0.5 most of the time.

% Supplementary Figure \ref{fig:conflict:exactly_one} shows the  . %todo
% Supplementary Figure \ref{fig:conflict:at_least_one}
% Supplementary Figure \ref{fig:conflict:exactly_two}.
% However, could be an artifact of cell-offspring cooperation.

% Interestingly, this strain appears to be using ``Is Child Cell Of,'' ``Is Parent Cell Of,'' ``Is Child Group Of,'' and the raw kin group ID of neighbors (but not the raw kin group ID of self).

% %todo change wording, not sure what you're hoping to say with first phrase
% To more conclusively, in future work we plan to look at the reproductive output at the interior of multicellular groups to see if these cells diminish their own reproductive output and to control for cellular relatedness when calculating this element of group cooperation.

% We analyzed other traits characteristic of multicellularity, as well.

% We did not see widespread resource sharing over the evolutionary history.
% Only tiny fractions of cells, less than 1\%, appeared to share tiny amounts of the resource, less than 1\% of the resource required to reproduce per update (Supplementary Figures \ref{fig:resource_sharing:fraction_sharing_monoculture} and \ref{fig:resource_sharing:sharing_amount_monoculture}).
% However, in ten strains sampled along the evolutionary lineage, context-dependent resource-sending behavior contributed to fitness (Supplementary Figure \ref{fig:writable_perturbation}).
% Interestingly, another strain present in the population that was not the subject of analysis appeared to evolve widespread, higher-throughput resource sharing after stint 60 (Supplementary Figures \ref{fig:resource_sharing:fraction_sharing_evolve} and \ref{fig:resource_sharing:sharing_amount_evolve}).

% We observed widespread apoptosis over the evolutionary history.
% As shown in Supplementary Figure \ref{fig:apoptosis:apoptosis_stint}, in most monocultured representative strains, between 10 and 30\% of cell deaths were due to apoptosis.
% However, it is not clear if this apoptosis returned a fitness benefit.
% No context-dependent apoptosis behavior contributed to fitness for any sampled strain over the evolutionary history (Supplementary Figure \ref{fig:writable_perturbation}).
% It is unclear whether context-independent apoptosis might have a fitness benefit.

% The behavior seems like cells from within a group don't explicitly recognize one another according to group membership. However, they do recognize their cell parent/child and arrange themselves in a somewhat linear fashion that minimizes the potential for conflict between cells.
% Overall, there are elements of cooperation but is unclear to what extent this particular strain constitutes a proper fraternal transition in individuality.


\subsection{Fitness Complexity}

Figure \ref{fig:complexity:genetic} plots critical fitness complexity of specimens drawn from across the case study's evolutionary history.

Critical fitness complexity reaches more than 20 under morph $b$, jumps to more than 40 under morph $d$, and drops to slightly more than 30 for morph $e$.
Critical fitness complexity reaches a peak of 48 sites around stint 39, then levels out and decreases.
This decrease may in part be due to declining sensitivity of competition experiments due to slower simulation resulting in execution of fewer updates within the fixed-duration jobs (Supplementary Figure \ref{fig:num_updates_elapsed_heatmap}).

Phylogenetic analysis (Supplementary Figure \ref{fig:phylo_parsimony_tree}) suggests independent origins of the critical fitness complexity in morph $d$ and morph $e$ --- the morph $d$ specimen from stint 14 is more closely related to the morph $b$ specimen from stint 13 than to the morph $e$ specimen from stint 15.
Likewise, specimens of lower complexity morphs $i$ and $b$ that appear past stint 70 appear to have independent evolutionary origins.

% The evolution of morph $g$ is not associated with a clear change in critical fitness complexity.

% Interpolated fitness complexity remains much more steady over the course of evolutionary history, although the certainty of our estimate of interpolated fitness complexity varies greatly.
% Figure \ref{fig:fitness_complexity:interpolated_fitness_complexity} shows this.
% Sometimes we were not able to estimate composite fitness complexity because the phenotype neutral nopout was not truly neutral --- it was significantly less fit than wildtype.

% Figure \ref{fig:fitness_complexity:composite_fitness_complexity} shows composite fitness complexity, the sum of interpolated fitness complexity and critical fitness complexity.
% The highest best estimate of composite fitness complexity was 64 sites at stint 61.
% The highest lower bound estimate of composite fitness complexity was 49 sites at stint 70.

\subsection{Interface Complexity}

Supplementary Figure \ref{fig:interface_complexity} summarizes cardinal processor interface complexity, as well as its constituent components, for specimens drawn from across the case study's evolutionary history.

Notably, cardinal processor interface complexity more than doubles from 6 interactions to 17 interactions coincident with the emergence of morph $e$ (Supplementary Figure \ref{fig:interface_complexity:cardinal_interface_complexity}).
This is due to simultaneous increases in extrospective state sensing (2 to 9 states; Supplementary Figure \ref{fig:interface_complexity:extrospective_interface_complexity}), introspective state sensing (1 to 4 states; Supplementary Figure \ref{fig:interface_complexity:introspective_interface_complexity}), and writable state usage (1 to 2 states; Supplementary Figure \ref{fig:interface_complexity:writable_interface_complexity}).

The emergence of morph $g$ coincided with an increase in writable state interface complexity from 1 to 3, as shown in Supplementary Figure \ref{fig:interface_complexity:writable_interface_complexity}.
However, morph $g$ was not associated with other changes in other aspects of cardinal processor interface complexity.
The greatest observed cardinal processor interface complexity was 22 interactions at stints 54 and 67.

\subsection{Genome Size}

Supplementary Figure \ref{fig:genome_size} shows evolutionary trajectories of three genome size metrics in sampled focal strain specimens.
Instruction count and module count increased from 100 and 5 to around 800 and 30, respectively, between stints 0 and 40.
Within this period at stint 24, instruction count jumped from around 600 to more than 800 and module count jumped by about 5.
This was coincident with detection in our adaptation assays of population-level sign-change mediation of adaptation by the background strain (Supplementary Figure \ref{fig:fitness_gain_or_loss}).
In sampled specimen fitness assays at stint 24, we detected significant increases in fitness in the presence of the background strain but no significant change in fitness in its absence.

Between stints 40 and 90, module count gradually increased to around 40 while instruction count remained stable.
Then, at stint 93, instruction count jumped to around 1,500 and module count jumped to around 60.
This was coincident with the strong fitness differentials observed at stint 93 (Supplementary Figure \ref{fig:median_fitness_differential_symlog}).

To better understand the functional effects of changes in genome size, we additionally measured the number of instructions that affected agent phenotype, shown in Figure \ref{fig:phenotype_complexity}.
This measure can be considered akin to a count of ``active'' sites.

Active site count varied greatly stint to stint.
The median value increased from nearly 0 to around 200 sites between stints 0 and 40.
Between stints 40 and 90, we observed active site values ranging from less than 100 to more than 500.
Morph $g$ specimens appear to show particularly great variance in active site count.
The first observed morph $g$ specimen at stint 45 exhibited relatively low active site count of around 100 active sites.
The highest active site count values of around 700 were measured from three specimens of morphs $e$ and $g$ in the last ten stints.
